\begin{theorem}[Hilbert Nullstellensatz]
\label{theorem-nullstellensatz}
Let $k$ be a field.
\begin{enumerate}
\item
\label{item-finite-kappa}
For any maximal ideal $\mathfrak m \subset k[x_1, \ldots, x_n]$
the field extension $\kappa(\mathfrak m)/k$ is finite.
\item
\label{item-polynomial-ring-Jacobson}
Any radical ideal $I \subset k[x_1, \ldots, x_n]$
is the intersection of maximal ideals containing it.
\end{enumerate}
The same is true in any finite type $k$-algebra.
\end{theorem}

\begin{proof}
It is enough to prove part (\ref{item-finite-kappa}) of
the theorem for the case of a polynomial
algebra $k[x_1, \ldots, x_n]$, because any finitely generated
$k$-algebra is a quotient of such a polynomial algebra.
We prove this by induction on $n$. The case $n = 0$ is clear.
Suppose that $\mathfrak m$ is a maximal ideal in $k[x_1, \ldots, x_n]$.
Let $\mathfrak p \subset k[x_n]$ be the intersection
of $\mathfrak m$ with $k[x_n]$.

\medskip\noindent
If $\mathfrak p \not = (0)$,
then $\mathfrak p$ is maximal and generated by an irreducible
monic polynomial $P$ (because of the Euclidean algorithm
in $k[x_n]$). Then
$k' = k[x_n]/\mathfrak p$ is a finite field extension of $k$
and contained in $\kappa(\mathfrak m)$. In this case
we get a surjection
$$
k'[x_1, \ldots, x_{n-1}]
\to
k'[x_1, \ldots, x_n] =
k' \otimes_k k[x_1, \ldots, x_n]
\longrightarrow
\kappa(\mathfrak m)
$$
and hence we see that $\kappa(\mathfrak m)$ is a finite
extension  of $k'$ by induction hypothesis. Thus $\kappa(\mathfrak m)$
is finite over $k$ as well.

\medskip\noindent
If $\mathfrak p = (0)$ we consider the ring
extension $k[x_n] \subset k[x_1, \ldots, x_n]/\mathfrak m$.
This is a finitely generated ring extension, hence
of finite presentation by
Lemmas \ref{lemma-obvious-Noetherian} and
\ref{lemma-Noetherian-finite-type-is-finite-presentation}.
Thus the image of $\Spec(k[x_1, \ldots, x_n]/\mathfrak m)$
in $\Spec(k[x_n])$ is constructible by
Theorem \ref{theorem-chevalley}. Since the image
contains $(0)$ we conclude that it contains a standard
open $D(f)$ for some $f\in k[x_n]$ nonzero. Since clearly
$D(f)$ is infinite we get a contradiction with the
assumption that $k[x_1, \ldots, x_n]/\mathfrak m$ is
a field (and hence has a spectrum consisting of one point).

\medskip\noindent
Proof of (\ref{item-polynomial-ring-Jacobson}). Let
$I \subset R$ be a radical ideal, with $R$ of finite type over $k$.
Let $f \in R$, $f \not \in I$. We have to find a maximal ideal
$\mathfrak m \subset R$ with $I \subset \mathfrak m$ and
$f \not \in \mathfrak m$. The ring $(R/I)_f$ is nonzero, since
$1 = 0$ in this ring would mean $f^n \in I$ and since $I$ is
radical this would mean $f \in I$ contrary to our assumption on $f$.
Thus we may choose a maximal ideal $\mathfrak m'$
in $(R/I)_f$, see Lemma \ref{lemma-Zariski-topology}.
Let $\mathfrak m \subset R$
be the inverse image of $\mathfrak m'$ in $R$. We see that
$I \subset \mathfrak m$
and $f \not \in \mathfrak m$. If we show that $\mathfrak m$ is a maximal
ideal of $R$, then we are done. We clearly have
$$
k \subset R/\mathfrak m \subset \kappa(\mathfrak m').
$$
By part (\ref{item-finite-kappa}) the field extension
$\kappa(\mathfrak m')/k$ is finite. Hence
$R/\mathfrak m$ is a field by Fields, Lemma
\ref{fields-lemma-subalgebra-algebraic-extension-field}.
Thus $\mathfrak m$ is maximal and the proof is complete.
\end{proof}